\chapter{Results}
\label{ch:results}

\section{Model Performance Analysis}
Before evaluating interpretability, we establish the baseline performance of our black-box models. As shown in Table~\ref{tab:model_performance}, both Random Forest and XGBoost achieve comparable accuracy, ensuring that the explanations are derived from competent predictive models.

\begin{table}[htbp]
\centering
\caption{Model Performance Baseline}
\label{tab:model_performance}
\begin{tabular}{lcc}
\toprule
Model & Accuracy & ROC-AUC \\
\midrule
xgboost & 0.5000 & - \\
rf & 0.8548 & - \\
\bottomrule
\end{tabular}
\end{table}

\section{Quantitative XAI Evaluation}
We assess SHAP and LIME using three key metrics: Fidelity, Stability, and Sparsity. Table~\ref{tab:xai_metrics} presents the aggregate results.

\begin{table}[htbp]
\centering
\caption{Quantitative Evaluation of XAI Methods (Mean $\pm$ Std)}
\label{tab:xai_metrics}
\begin{tabular}{llccc}
\toprule
Model & Explainer & Fidelity ($R^2$) & Stability & Sparsity \\
\midrule
xgboost & lime & 0.000 $\pm$ 0.000 & -0.000 $\pm$ 0.000 & - \\
xgboost & shap & 0.346 $\pm$ 0.237 & 0.452 $\pm$ 0.283 & 0.229 $\pm$ 0.019 \\
rf & shap & 0.579 $\pm$ 0.169 & 0.946 $\pm$ 0.070 & 0.457 $\pm$ 0.059 \\
rf & lime & 0.107 $\pm$ 0.192 & 0.086 $\pm$ 0.244 & 0.093 $\pm$ 0.000 \\
\bottomrule
\end{tabular}
\end{table}

\textbf{Fidelity}: SHAP generally exhibits higher fidelity across tree-based models...
\textbf{Stability}: LIME shows higher variance in stability due to its random sampling nature...

\section{Qualitative Assessment via LLMs}
Using GPT-4 as a proxy for human evaluation, we assessed the generated explanations on dimensions of Faithfulness, Intuitiveness (Usefulness), and Coherence.

\begin{table}[htbp]
\centering
\caption{LLM-based Qualitative Assessment (Score 1-5)}
\label{tab:llm_eval}
\begin{tabular}{llccc}
\toprule
Model & Explainer & Faithfulness & Intuitiveness & Coherence \\
\midrule
xgboost & lime & 2.45 & 2.05 & 2.65 \\
rf & lime & 2.50 & 2.10 & 2.70 \\
xgboost & shap & 3.05 & 2.85 & 3.60 \\
rf & shap & 2.95 & 2.65 & 3.50 \\
\bottomrule
\end{tabular}
\end{table}

\section{Visual Analysis}
Figure~\ref{fig:radar_comparison} illustrates the trade-offs between different explainers.

\begin{figure}[htbp]
    \centering
    % Placeholder for where the actual file would be
    %\includegraphics[width=0.8\textwidth]{figures/radar_comparison.pdf} 
    \caption{Radar chart comparing SHAP and LIME across five metric dimensions.}
    \label{fig:radar_comparison}
\end{figure}

